\chapter{Полное корреляционное ядро и возврат конечной локальности}
\label{ch:v8-full-correlation-kernel-locality-reconstruction-gate}

\section{Пропущенный промежуточный объект}

Ранний TOE--UGSM-мост различал фундаментальный нелокальный оператор
$\widehat C$ и его спектрально-моментную редукцию. Однако обратно-спектральный
гейт Тома V проверял скалярные спектральные данные, а в качестве
положительной крайности оставлял уже полный спектральный тройной набор.
Между этими крайностями существует более слабый объект:
\begin{equation}
 (\mathcal A_{\rm obs},C_\tau),
 \qquad C_\tau=e^{-\tau H},
 \label{eq:v8-full-kernel-object}
\end{equation}
где $\mathcal A_{\rm obs}$ фиксирует различимые минимальные наблюдаемые, а
$C_\tau$ сохраняется как полный оператор, а не только как список
собственных значений или след.

\section{Повторный тест на старой коспектральной паре}

Используем точный контрпример Тома V: связную звезду $K_{1,4}$ и
$C_4\sqcup\{\mathrm{pt}\}$. Их матрицы смежности $A_1,A_2$ имеют один
характеристический многочлен, а положительные генераторы
\begin{equation}
 H_i=3I-A_i
\end{equation}
имеют общий спектр $(1,3,3,3,5)$. Поэтому для всякого $\tau>0$
совпадают тепловые следы и вообще все скалярные спектральные функционалы.

Полные ядра в выделенной диагональной алгебре $\mathbb C^5$, напротив,
различны даже после перебора всех перестановок пяти минимальных проекторов.
Для $\tau=0.1,0.3,1,2$ минимальные фробениусовы расстояния равны
соответственно приблизительно
\begin{equation}
 0.2011,\quad 0.3239,\quad 0.2206,\quad 0.08815.
\end{equation}
Следовательно, различие связности было уничтожено не тепловым оператором,
а последующим переходом к одному спектру или следу.

\section{Точный возврат генератора}

Поскольку $C_\tau$ положительно определён, функциональное исчисление даёт
\begin{equation}
 H=-\tau^{-1}\log C_\tau.
 \label{eq:v8-kernel-log-return}
\end{equation}
В конечной модели внедиагональная отрицательная опора $H$ точно возвращает
матрицу смежности. Численная ошибка реконструкции генератора во всех тестах
не превосходит $2\cdot10^{-13}$. Для восстановления одной лишь опоры не
нужно знать ни скалярный сдвиг $3I$, ни точное положительное значение
$\tau$: ошибка времени только положительно перемасштабирует $H$.

Этот возврат ковариантен при одновременной перестановке минимальных
проекторов $\mathbb C^5$. Тем самым в тест не внесены имена вершин, но
внесено более существенное исходное данное --- сама алгебра различимых
событий.

\begin{theorem}[Уточнение обратно-спектральный границы]
Спектр $C_\tau$ и все его скалярные следовые функционалы не определяют
конечную локальность. Полный положительный полугрупповой оператор
$C_\tau$ в фиксированной коммутативной алгебре наблюдаемых определяет свой
генератор функциональным логарифмом и в данной паре точно различает
связную и несвязную геометрии. Поэтому отрицательный результат Тома V
сохраняется для слабого спектрального объекта, но не распространяется на
полное двухточечное ядро с алгеброй наблюдаемых.
\end{theorem}

\section{Почему это ещё не основание Тома VIII}

В конечном аудите диагональная алгебра $\mathbb C^5$ задана заранее. Не
выведено, какая алгебра наблюдаемых канонична на физическом носителе S2T;
не доказаны марковская нормировка, сильная локальность и происхождение
семейства $C_\tau$ из одного действия. Тем более пока не получен
нефакторизованный оператор, связывающий рёберный, концевой и семейный
носители.

\begin{nogocriterion}[Запрет снова сжимать источник до спектра]
Следующий тест не имеет права заменять $C_\tau$ его собственными значениями,
тепловым следом или конечным набором моментов до проверки локальности и
межсекторных матричных элементов. Но также нельзя считать произвольно
выбранный базис физической геометрией: алгебра наблюдаемых и полугруппа
должны быть выведены из накопленной структуры.
\end{nogocriterion}

\begin{protocol}
\begin{enumerate}
 \item старый косpectral-контрпример: \textbf{сохранён};
 \item совпадение всех скалярных спектральных функционалов: \textbf{да};
 \item различимость полных ядер: \textbf{да};
 \item возврат генератора и конечной связности: \textbf{точный};
 \item полный спектральный тройной набор предполагался: \textbf{нет};
 \item алгебра наблюдаемых выведена: \textbf{нет};
 \item физическая S2T-полугруппа выведена: \textbf{нет};
 \item промежуточный полного ядра маршрут: \textbf{переоткрыт}.
\end{enumerate}
\end{protocol}

Машинный сертификат сохранён в
\path{s2t_v8_full_correlation_kernel_locality_reconstruction_gate_results.json}.